For few-shot learning, we evaluate each example in the evaluation set by randomly drawing $K$ examples from that task’s training set as conditioning, delimited by 1 or 2 newlines depending on the task.  For LAMBADA and Storycloze there is no supervised training set available so we draw conditioning examples from the development set and evaluate on the test set.  For Winograd (the original, not SuperGLUE version) there is only one dataset, so we draw conditioning examples directly from it.

$K$ can be any value from 0 to the maximum amount allowed by the model’s context window, which is $n_{\mathrm{ctx}}=2048$ for all models and typically fits $10$ to $100$ examples.  Larger values of $K$ are usually but not always better, so when a separate development and test set are available, we experiment with a few values of $K$ on the development set and then run the best value on the test set.  For some tasks (see Appendix \ref{appendix:task_phrasing}) we also use a natural language prompt in addition to (or for $K=0$, instead of) demonstrations.

On tasks that involve choosing one correct completion from several options (multiple choice), we provide $K$ examples of context plus correct completion, followed by one example of context only, and compare the LM likelihood of each completion.  For most tasks we compare the per-token likelihood (to normalize for length), however on a small number of datasets (ARC, OpenBookQA, and RACE) we gain additional benefit as measured on the development set by normalizing by the unconditional probability of each completion, by computing $\frac{P(\mathrm{completion} | \mathrm{context})}{P(\mathrm{completion} | \mathrm{answer\_context})}$, where $\mathrm{answer\_context}$ is the string \texttt{"Answer: "} or \texttt{"A: "} and is used to prompt that the completion should be an answer but is otherwise generic.

On tasks that involve binary classification, we give the options more semantically meaningful names (e.g. ``True" or ``False" rather than 0 or 1) and then treat the task like multiple choice; we also sometimes frame the task similar to what is done by \cite{raffel2019t5} (see Appendix \ref{appendix:task_phrasing}) for details.

On tasks with free-form completion, we use beam search with the same parameters as \cite{raffel2019t5}: a beam width of 4 and a length penalty of $\alpha = 0.6$.  We score the model using F1 similarity score, BLEU, or exact match, depending on what is standard for the dataset at hand.

Final results are reported on the test set when publicly available, for each model size and learning setting (zero-, one-, and few-shot).  When the test set is private, our model is often too large to fit on the test server, so we report results on the development set.  We do submit to the test server on a small number of datasets (SuperGLUE, TriviaQA, PiQa) where we were able to make submission work, and we submit only the 200B few-shot results, and report development set results for everything else. 
